\sectionheading{5.37 Conditional Logic Statement}

Program-এ সব Statement সবসময় Execute করতে হয় না।
কোনো Condition সত্য বা মিথ্যা হওয়ার উপর নির্ভর করে
নির্দিষ্ট Code Block Execute করার জন্য
\textbf{Conditional Logic Statement} ব্যবহার করা হয়।

%================================================
% Definition
%================================================

\booksubheading{Conditional Statement কী?}

\begin{definitionbox}

কোনো Condition পরীক্ষা করে তার ফলাফল অনুযায়ী Program-এর
নির্দিষ্ট Statement বা Code Block Execute করার পদ্ধতিকে
\textbf{Conditional Statement} বলে।

\end{definitionbox}

\begin{center}

\textbf{Condition পরীক্ষা}
\(\longrightarrow\)
\textenglish{\ttfamily True / False}
\(\longrightarrow\)
\textbf{প্রয়োজনীয় Code Block Execute}

\end{center}

উদাহরণ—

\begin{codebox}
\begin{english}
\ttfamily
age = 18\\
if age >= 18:\\
\hspace*{4mm}print("Adult")
\end{english}
\end{codebox}

এখানে
\textenglish{\ttfamily age >= 18}
Condition সত্য হলে
\textenglish{\ttfamily Adult}
প্রদর্শিত হবে।

%================================================
% Main Forms
%================================================

\booksubheading{Python-এর প্রধান Conditional Statement}

\begin{tabularx}
{\linewidth}
{
|>{\bfseries\RaggedRight\arraybackslash}p{0.31\linewidth}
|Y|
}

\hline
\rowcolor{softBlue}

\multicolumn{1}{|c|}{\textbf{Conditional Form}}
&
\multicolumn{1}{c|}{\textbf{ব্যবহার}}
\\
\hline

\textenglish{\ttfamily if}
&
Condition True হলে Code Execute।
\\
\hline

\textenglish{\ttfamily if ... else}
&
True/False অনুযায়ী দুইটি Block-এর একটি Execute।
\\
\hline

\textenglish{\ttfamily if ... elif ... else}
&
একাধিক Condition পর্যায়ক্রমে পরীক্ষা।
\\
\hline

Nested \textenglish{\ttfamily if}
&
একটি
\textenglish{\ttfamily if}
-এর ভিতরে আরেকটি
\textenglish{\ttfamily if}।
\\
\hline

\end{tabularx}

%================================================
% IF
%================================================

\booksubheading{১. \textenglish{\ttfamily if} Statement}

কোনো Condition
\textenglish{\ttfamily True}
হলে নির্দিষ্ট Code Block Execute করতে
\textenglish{\ttfamily if}
Statement ব্যবহার করা হয়।

\textbf{Syntax—}

\begin{codebox}
\begin{english}
\ttfamily
if condition:\\
\hspace*{4mm}statement
\end{english}
\end{codebox}

\begin{examplebox}{উদাহরণ}

\begin{codebox}
\begin{english}
\ttfamily
marks = 80\\
if marks >= 40:\\
\hspace*{4mm}print("Pass")
\end{english}
\end{codebox}

\textbf{Output:}
\textenglish{\ttfamily Pass}

\end{examplebox}

\begin{center}

\textbf{Condition True}
\(\longrightarrow\)
\textenglish{\ttfamily if}
Block Execute

\qquad

\textbf{Condition False}
\(\longrightarrow\)
Block Skip

\end{center}

%================================================
% Colon + indentation
%================================================

\booksubheading{Colon \textenglish{\ttfamily :} ও Indentation}

Python Conditional Statement-এ দুটি বিষয় খুব গুরুত্বপূর্ণ—

\begin{itemize}

\item Condition-এর শেষে
\textenglish{\ttfamily :}
দিতে হয়।

\item
\textenglish{\ttfamily if}
Block-এর Statement Indented হতে হয়।

\item সাধারণভাবে 4 Spaces Indentation ব্যবহার করা হয়।

\end{itemize}

\noindent
\begin{minipage}[t]{0.49\linewidth}

\begin{sideexamplebox}{ভুল}

\begin{sidecodebox}
\begin{english}
\ttfamily
if marks >= 40\\
print("Pass")
\end{english}
\end{sidecodebox}

Colon নেই এবং Indentation-ও নেই।

\end{sideexamplebox}

\end{minipage}
\hfill
\begin{minipage}[t]{0.49\linewidth}

\begin{sideexamplebox}{সঠিক}

\begin{sidecodebox}
\begin{english}
\ttfamily
if marks >= 40:\\
\hspace*{4mm}print("Pass")
\end{english}
\end{sidecodebox}

Colon ও Indentation সঠিক।

\end{sideexamplebox}

\end{minipage}

\begin{rememberbox}

\textbf{\color{deepOrange}MCQ Trap}

Python-এ Code Block নির্ধারণে
\textbf{Indentation} গুরুত্বপূর্ণ।

\begin{center}

\textbf{C/C++}
\(\longrightarrow\)
\textenglish{\ttfamily \{ \}}

\qquad

\textbf{Python}
\(\longrightarrow\)
\textbf{Indentation}

\end{center}

\end{rememberbox}

%================================================
% IF ELSE
%================================================

\booksubheading{২. \textenglish{\ttfamily if ... else} Statement}

Condition True হলে একটি Block এবং False হলে অন্য Block
Execute করার জন্য
\textenglish{\ttfamily if ... else}
ব্যবহার করা হয়।

\textbf{Syntax—}

\begin{codebox}
\begin{english}
\ttfamily
if condition:\\
\hspace*{4mm}statement1\\
else:\\
\hspace*{4mm}statement2
\end{english}
\end{codebox}

\noindent
\begin{minipage}[t]{0.49\linewidth}

\begin{sideexamplebox}{Pass–Fail}

\begin{sidecodebox}
\begin{english}
\ttfamily
marks = 35\\
if marks >= 40:\\
\hspace*{4mm}print("Pass")\\
else:\\
\hspace*{4mm}print("Fail")
\end{english}
\end{sidecodebox}

Output:
\textenglish{\ttfamily Fail}

\end{sideexamplebox}

\end{minipage}
\hfill
\begin{minipage}[t]{0.49\linewidth}

\begin{sideexamplebox}{Even–Odd}

\begin{sidecodebox}
\begin{english}
\ttfamily
n = int(input())\\
if n \% 2 == 0:\\
\hspace*{4mm}print("Even")\\
else:\\
\hspace*{4mm}print("Odd")
\end{english}
\end{sidecodebox}

\textenglish{\ttfamily n \% 2 == 0}
হলে Even, অন্যথায় Odd।

\end{sideexamplebox}

\end{minipage}

\begin{rememberbox}

\textbf{\color{deepOrange}মনে রাখবে}

\begin{center}

\textenglish{\ttfamily n \% 2 == 0}
\(\longrightarrow\)
Even

\qquad

\textenglish{\ttfamily n \% 2 != 0}
\(\longrightarrow\)
Odd

\end{center}

\textenglish{\ttfamily if ... else}
-এ দুইটি Block-এর মধ্যে একটিমাত্র Execute হয়।

\end{rememberbox}

%================================================
% ELIF
%================================================

\booksubheading{৩. \textenglish{\ttfamily if ... elif ... else}}

একাধিক Condition পর্যায়ক্রমে পরীক্ষা করার জন্য
\textenglish{\ttfamily elif}
ব্যবহার করা হয়।

\begin{center}
\textenglish{\ttfamily elif}
=
\textbf{else if}
\end{center}

\textbf{Syntax—}

\begin{codebox}
\begin{english}
\ttfamily
if condition1:\\
\hspace*{4mm}statement1\\
elif condition2:\\
\hspace*{4mm}statement2\\
else:\\
\hspace*{4mm}statement3
\end{english}
\end{codebox}

\begin{examplebox}{Grade নির্ণয়}

\begin{codebox}
\begin{english}
\ttfamily
marks = 75\\
if marks >= 80:\\
\hspace*{4mm}print("A+")\\
elif marks >= 70:\\
\hspace*{4mm}print("A")\\
elif marks >= 60:\\
\hspace*{4mm}print("A-")\\
else:\\
\hspace*{4mm}print("Below A-")
\end{english}
\end{codebox}

\textbf{Output:}
\textenglish{\ttfamily A}

\end{examplebox}

%================================================
% How elif works
%================================================

\booksubheading{\textenglish{\ttfamily elif} কীভাবে কাজ করে?}

Python ওপর থেকে নিচে Condition পরীক্ষা করে।
যে Condition প্রথম
\textenglish{\ttfamily True}
হয়, তার Block Execute হয় এবং পরবর্তী
\textenglish{\ttfamily elif/else}
Blockগুলো আর Execute হয় না।

\begin{center}

\textbf{Condition 1?}
\(\longrightarrow\)
False
\(\longrightarrow\)
\textbf{Condition 2?}
\(\longrightarrow\)
False
\(\longrightarrow\)
\textbf{Condition 3?}
\(\longrightarrow\)
\textbf{else}

\end{center}

\begin{rememberbox}

\textbf{\color{deepOrange}গুরুত্বপূর্ণ}

একাধিক Condition True হলেও
\textbf{প্রথম True Condition-এর Block-ই Execute হয়}।

\end{rememberbox}

%================================================
% Condition Order
%================================================

\booksubheading{Condition-এর Order গুরুত্বপূর্ণ}

\noindent
\begin{minipage}[t]{0.49\linewidth}

\begin{sideexamplebox}{ভুল Order}

\begin{sidecodebox}
\begin{english}
\ttfamily
marks = 85\\
if marks >= 40:\\
\hspace*{4mm}print("Pass")\\
elif marks >= 80:\\
\hspace*{4mm}print("A+")
\end{english}
\end{sidecodebox}

Output:
\textenglish{\ttfamily Pass}

\end{sideexamplebox}

\end{minipage}
\hfill
\begin{minipage}[t]{0.49\linewidth}

\begin{sideexamplebox}{ভালো Order}

\begin{sidecodebox}
\begin{english}
\ttfamily
marks = 85\\
if marks >= 80:\\
\hspace*{4mm}print("A+")\\
elif marks >= 40:\\
\hspace*{4mm}print("Pass")
\end{english}
\end{sidecodebox}

Output:
\textenglish{\ttfamily A+}

\end{sideexamplebox}

\end{minipage}

\begin{rememberbox}

Range বা Grade পরীক্ষা করার সময় সাধারণত
\textbf{বড়/Specific Condition আগে}
দেওয়া সুবিধাজনক।

\end{rememberbox}

%================================================
% Optional elif/else
%================================================

\booksubheading{\textenglish{\ttfamily elif} ও \textenglish{\ttfamily else} কি বাধ্যতামূলক?}

না। প্রয়োজন অনুযায়ী শুধু
\textenglish{\ttfamily if}
ব্যবহার করা যায়।

\begin{codebox}
\begin{english}
\ttfamily
if x > 0:\\
\hspace*{4mm}print("Positive")
\end{english}
\end{codebox}

আবার
\textenglish{\ttfamily else}
ছাড়াও
\textenglish{\ttfamily elif}
ব্যবহার করা যায়।

\begin{codebox}
\begin{english}
\ttfamily
if x > 0:\\
\hspace*{4mm}print("Positive")\\
elif x < 0:\\
\hspace*{4mm}print("Negative")
\end{english}
\end{codebox}

\begin{rememberbox}

\textbf{\color{deepOrange}MCQ তথ্য}

\begin{center}

\textenglish{\ttfamily if}
\(\longrightarrow\)
প্রধান Condition

\qquad

\textenglish{\ttfamily elif}
\(\longrightarrow\)
0 বা একাধিক হতে পারে

\qquad

\textenglish{\ttfamily else}
\(\longrightarrow\)
Optional

\end{center}

\end{rememberbox}

%================================================
% Nested if
%================================================

\booksubheading{৪. Nested \textenglish{\ttfamily if}}

একটি
\textenglish{\ttfamily if}
বা
\textenglish{\ttfamily else}
Block-এর ভিতরে আরেকটি
\textenglish{\ttfamily if}
Statement থাকলে তাকে
\textbf{Nested if} বলে।

\noindent
\begin{minipage}[t]{0.49\linewidth}

\begin{sideexamplebox}{Nested \textenglish{\ttfamily if}}

\begin{sidecodebox}
\begin{english}
\ttfamily
age = 20\\
citizen = True\\
if age >= 18:\\
\hspace*{4mm}if citizen:\\
\hspace*{8mm}print("Eligible")
\end{english}
\end{sidecodebox}

Outer Condition True হলে Inner Condition পরীক্ষা হয়।

\end{sideexamplebox}

\end{minipage}
\hfill
\begin{minipage}[t]{0.49\linewidth}

\begin{sideexamplebox}{Logical \textenglish{\ttfamily and} দিয়ে}

\begin{sidecodebox}
\begin{english}
\ttfamily
age = 20\\
citizen = True\\
if age >= 18 and citizen:\\
\hspace*{4mm}print("Eligible")
\end{english}
\end{sidecodebox}

অনেক ক্ষেত্রে Nested
\textenglish{\ttfamily if}
-এর Logic সংক্ষেপে লেখা যায়।

\end{sideexamplebox}

\end{minipage}

%================================================
% Comparison and Logical Operators
%================================================

\booksubheading{Condition-এ Comparison ও Logical Operator}

Conditional Statement-এ Comparison Operator খুব বেশি ব্যবহৃত হয়।

\begin{tabularx}
{\linewidth}
{
|>{\bfseries\centering\arraybackslash}p{0.15\linewidth}
|Y
|>{\bfseries\centering\arraybackslash}p{0.15\linewidth}
|Y|
}

\hline
\rowcolor{softBlue}

\textbf{Operator}
&
\textbf{অর্থ}
&
\textbf{Operator}
&
\textbf{অর্থ}
\\
\hline

\textenglish{\ttfamily ==}
&
সমান
&
\textenglish{\ttfamily !=}
&
অসমান
\\
\hline

\textenglish{\ttfamily >}
&
বড়
&
\textenglish{\ttfamily <}
&
ছোট
\\
\hline

\textenglish{\ttfamily >=}
&
বড় বা সমান
&
\textenglish{\ttfamily <=}
&
ছোট বা সমান
\\
\hline

\end{tabularx}

একাধিক Condition যুক্ত করতে
\textenglish{\ttfamily and},
\textenglish{\ttfamily or}
ও
\textenglish{\ttfamily not}
ব্যবহার করা হয়।

\noindent
\begin{minipage}[t]{0.32\linewidth}

\begin{sideexamplebox}{\textenglish{\ttfamily and}}

\begin{sidecodebox}
\begin{english}
\ttfamily
if age >= 18\\
\hspace*{3mm}and age <= 30:\\
\hspace*{4mm}print("Eligible")
\end{english}
\end{sidecodebox}

সব Condition True হতে হবে।

\end{sideexamplebox}

\end{minipage}
\hfill
\begin{minipage}[t]{0.32\linewidth}

\begin{sideexamplebox}{\textenglish{\ttfamily or}}

\begin{sidecodebox}
\begin{english}
\ttfamily
if day == "Friday"\\
\hspace*{3mm}or day == "Saturday":\\
\hspace*{4mm}print("Holiday")
\end{english}
\end{sidecodebox}

অন্তত একটি True হলেই যথেষ্ট।

\end{sideexamplebox}

\end{minipage}
\hfill
\begin{minipage}[t]{0.32\linewidth}

\begin{sideexamplebox}{\textenglish{\ttfamily not}}

\begin{sidecodebox}
\begin{english}
\ttfamily
if not raining:\\
\hspace*{4mm}print("Go outside")
\end{english}
\end{sidecodebox}

Logical Result বিপরীত করে।

\end{sideexamplebox}

\end{minipage}

%================================================
% Boolean directly
%================================================

\booksubheading{Boolean Value সরাসরি Condition হিসেবে}

Boolean Variable সরাসরি
\textenglish{\ttfamily if}
Statement-এ ব্যবহার করা যায়।

\begin{codebox}
\begin{english}
\ttfamily
logged\_in = True\\
if logged\_in:\\
\hspace*{4mm}print("Welcome")
\end{english}
\end{codebox}

\textbf{Output:}
\textenglish{\ttfamily Welcome}

\begin{rememberbox}

\textenglish{\ttfamily if logged\_in:}
লেখাই যথেষ্ট।

\textenglish{\ttfamily logged\_in == True}
লেখা সম্ভব হলেও সাধারণত প্রয়োজন নেই।

\end{rememberbox}

%================================================
% Truthy and Falsy
%================================================

\booksubheading{Truthy ও Falsy-এর প্রাথমিক ধারণা}

Python-এ Condition-এর মধ্যে কিছু Value
\textenglish{\ttfamily False}
হিসেবে বিবেচিত হয়। এগুলোকে সাধারণভাবে
\textbf{Falsy Value} বলা হয়।

\begin{tabularx}
{\linewidth}
{
|>{\bfseries\RaggedRight\arraybackslash}p{0.27\linewidth}
|Y|
}

\hline
\rowcolor{softBlue}

\multicolumn{1}{|c|}{\textbf{Falsy Type}}
&
\multicolumn{1}{c|}{\textbf{উদাহরণ}}
\\
\hline

Boolean
&
\textenglish{\ttfamily False}
\\
\hline

None
&
\textenglish{\ttfamily None}
\\
\hline

Zero
&
\textenglish{\ttfamily 0,\;0.0}
\\
\hline

Empty String
&
\textenglish{\ttfamily ""}
\\
\hline

Empty List / Tuple
&
\textenglish{\ttfamily [],\;()}
\\
\hline

Empty Set / Dictionary
&
\textenglish{\ttfamily set(),\;\{\}}
\\
\hline

\end{tabularx}

উদাহরণ—

\begin{codebox}
\begin{english}
\ttfamily
name = ""\\
if name:\\
\hspace*{4mm}print("Name exists")\\
else:\\
\hspace*{4mm}print("Empty")
\end{english}
\end{codebox}

\textbf{Output:}
\textenglish{\ttfamily Empty}

\begin{rememberbox}

\textbf{\color{deepOrange}MCQ তথ্য:}
\quad
\textenglish{\ttfamily 0},
Empty String এবং Empty Collection সাধারণত Falsy।

\end{rememberbox}

%================================================
% Chained comparison
%================================================

\booksubheading{Chained Comparison Condition-এ}

Python-এর Chained Comparison
\textenglish{\ttfamily if}
Statement-এ ব্যবহার করা যায়।

\begin{codebox}
\begin{english}
\ttfamily
age = 20\\
if 18 <= age <= 30:\\
\hspace*{4mm}print("Eligible")
\end{english}
\end{codebox}

এটি মূলত—

\begin{center}

\textenglish{\ttfamily age >= 18 and age <= 30}

\end{center}

-এর সমতুল্য Logic প্রকাশ করে।

%================================================
% Short if + conditional expression
%================================================

\booksubheading{Short-hand \textenglish{\ttfamily if} ও Conditional Expression}

\noindent
\begin{minipage}[t]{0.49\linewidth}

\begin{sideexamplebox}{Short-hand \textenglish{\ttfamily if}}

এক লাইনে ছোট
\textenglish{\ttfamily if}
লেখা সম্ভব—

\begin{sidecodebox}
\begin{english}
\ttfamily
if age >= 18: print("Adult")
\end{english}
\end{sidecodebox}

ছোট Logic-এর জন্য ব্যবহার করা যায়; বড় Program-এ Indented
Style বেশি readable।

\end{sideexamplebox}

\end{minipage}
\hfill
\begin{minipage}[t]{0.49\linewidth}

\begin{sideexamplebox}{Conditional Expression}

\begin{sidecodebox}
\begin{english}
\ttfamily
result = "Pass" if marks >= 40 else "Fail"
\end{english}
\end{sidecodebox}

সাধারণ Structure—

\begin{center}
\textbf{True Value}
\textenglish{\ttfamily if}
\textbf{Condition}
\textenglish{\ttfamily else}
\textbf{False Value}
\end{center}

এটিকে Conditional/Ternary Expression বলা হয়।

\end{sideexamplebox}

\end{minipage}

%================================================
% Comparison table
%================================================

\booksubheading{\textenglish{\ttfamily if}, \textenglish{\ttfamily elif}, \textenglish{\ttfamily else}-এর পার্থক্য}

\begin{tabularx}
{\linewidth}
{
|>{\bfseries\RaggedRight\arraybackslash}p{0.24\linewidth}
|Y|
}

\hline
\rowcolor{softBlue}

\multicolumn{1}{|c|}{\textbf{Keyword}}
&
\multicolumn{1}{c|}{\textbf{কাজ}}
\\
\hline

\textenglish{\ttfamily if}
&
প্রথম Condition পরীক্ষা করে।
\\
\hline

\textenglish{\ttfamily elif}
&
আগের Condition False হলে নতুন Condition পরীক্ষা করে।
\\
\hline

\textenglish{\ttfamily else}
&
আগের সব Condition False হলে Execute হয়।
\\
\hline

Nested \textenglish{\ttfamily if}
&
একটি
\textenglish{\ttfamily if}
Block-এর ভিতরে আরেক Condition পরীক্ষা করে।
\\
\hline

\end{tabularx}

%================================================
% Largest of three
%================================================

\booksubheading{তিন সংখ্যার মধ্যে বড় নির্ণয়ের উদাহরণ}

\begin{examplebox}{Program}

\begin{codebox}
\begin{english}
\ttfamily
a = 10\\
b = 25\\
c = 15\\[0.5mm]
if a >= b and a >= c:\\
\hspace*{4mm}print("a is largest")\\
elif b >= a and b >= c:\\
\hspace*{4mm}print("b is largest")\\
else:\\
\hspace*{4mm}print("c is largest")
\end{english}
\end{codebox}

\textbf{Output:}
\textenglish{\ttfamily b is largest}

\end{examplebox}

%================================================
% Common Mistakes
%================================================

\booksubheading{সাধারণ ভুল}

\begin{rememberbox}

\begin{itemize}

\item Condition-এর শেষে
\textenglish{\ttfamily :}
না দেওয়া।

\item Indentation না দেওয়া বা অসমান Indentation দেওয়া।

\item Comparison-এর জায়গায়
\textenglish{\ttfamily =}
ব্যবহার করা।

\item
\textenglish{\ttfamily ==}
এবং
\textenglish{\ttfamily =}
গুলিয়ে ফেলা।

\item
\textenglish{\ttfamily elif}
-এর পরিবর্তে
\textenglish{\ttfamily else if}
লেখা।

\item
\textenglish{\ttfamily True/False}
-এর Capital Letter ভুল করা।

\item Grade বা Range Condition ভুল Order-এ লেখা।

\item Nested
\textenglish{\ttfamily if}
-এর Indentation ভুল করা।

\item একাধিক Condition যুক্ত করতে
\textenglish{\ttfamily and/or}
ভুলভাবে ব্যবহার করা।

\item প্রয়োজন থাকা সত্ত্বেও সব Condition False হলে কী হবে তা
Handle না করা।

\end{itemize}

\end{rememberbox}

%================================================
% Final Reminder
%================================================

\begin{rememberbox}

\textbf{\color{deepOrange}মনে রাখবে}

\begin{tabularx}
{\linewidth}
{
>{\bfseries\RaggedRight\arraybackslash}p{0.24\linewidth}
Y
>{\bfseries\RaggedRight\arraybackslash}p{0.24\linewidth}
Y
}

\textenglish{\ttfamily if}
&
True হলে Execute
&
\textenglish{\ttfamily else}
&
False হলে Execute
\\[0.7mm]

\textenglish{\ttfamily elif}
&
একাধিক Condition পরীক্ষা
&
Nested \textenglish{\ttfamily if}
&
\textenglish{\ttfamily if}-এর ভিতরে \textenglish{\ttfamily if}
\\[0.7mm]

Colon
&
Condition শেষে
\textenglish{\ttfamily :}
&
Block
&
Indentation
\\[0.7mm]

\textenglish{\ttfamily =}
&
Assignment
&
\textenglish{\ttfamily ==}
&
Comparison
\\[0.7mm]

\textenglish{\ttfamily and}
&
সব Condition True
&
\textenglish{\ttfamily or}
&
অন্তত একটি True
\\[0.7mm]

\textenglish{\ttfamily not}
&
Logical Result বিপরীত
&
\textenglish{\ttfamily elif}
&
else if
\\

\end{tabularx}

\medskip

\begin{center}

\textbf{Condition}
\(\longrightarrow\)
\textenglish{\ttfamily True}
\(\longrightarrow\)
\textbf{একটি পথ}

\qquad

\textbf{Condition}
\(\longrightarrow\)
\textenglish{\ttfamily False}
\(\longrightarrow\)
\textbf{অন্য পথ}

\medskip

\textbf{Python Block}
\(\longrightarrow\)
\textbf{Indentation}

\qquad

\textenglish{\ttfamily n \% 2 == 0}
\(\longrightarrow\)
\textbf{Even}

\end{center}

\end{rememberbox}

\small

%================================================
% 5.38 Loop
%================================================

